\documentclass[12pt, letterpaper]{article}
\usepackage[margin=0.75in, top=0.75in, bottom=0.75in]{geometry}
\usepackage{booktabs}
\usepackage{longtable}
\usepackage{array}
\usepackage{xcolor}
\usepackage{colortbl}
\usepackage{hyperref}
\usepackage{helvet}
\renewcommand{\familydefault}{\sfdefault}

% Color definitions
\definecolor{headerbg}{RGB}{31, 73, 125}
\definecolor{headertext}{RGB}{255, 255, 255}
\definecolor{row1}{RGB}{235, 241, 250}
\definecolor{row2}{RGB}{255, 255, 255}
\definecolor{domainA}{RGB}{226, 239, 218}
\definecolor{domainB}{RGB}{226, 239, 218}
\definecolor{domainC}{RGB}{234, 241, 251}
\definecolor{domainD}{RGB}{234, 241, 251}
\definecolor{domainE}{RGB}{252, 228, 214}
\definecolor{domainF}{RGB}{255, 242, 204}
\definecolor{domainG}{RGB}{255, 242, 204}
\definecolor{domainH}{RGB}{243, 230, 245}
\definecolor{domainI}{RGB}{235, 241, 250}

\hypersetup{
    colorlinks=true,
    linkcolor=blue,
    urlcolor=blue
}

\begin{document}

\begin{center}
    {\LARGE \textbf{8th Grade Math --- Full Year Curriculum}}\\[4pt]
    {\large Chicago Public Schools (CPS Skyline) $\cdot$ 25 Lessons across 9 Units}
\end{center}

\vspace{6pt}

\small
\begin{longtable}{
    >{\centering\arraybackslash}p{0.4cm}   % #
    >{\raggedright\arraybackslash}p{3.2cm}  % Unit
    >{\raggedright\arraybackslash}p{3.0cm}  % Lesson Title
    >{\raggedright\arraybackslash}p{5.6cm}  % Description
    >{\centering\arraybackslash}p{2.2cm}    % Standard
}

% Header
\rowcolor{headerbg}
\textcolor{headertext}{\textbf{\#}} &
\textcolor{headertext}{\textbf{Unit}} &
\textcolor{headertext}{\textbf{Lesson Title}} &
\textcolor{headertext}{\textbf{Description}} &
\textcolor{headertext}{\textbf{Standard}} \\
\endfirsthead

\rowcolor{headerbg}
\textcolor{headertext}{\textbf{\#}} &
\textcolor{headertext}{\textbf{Unit}} &
\textcolor{headertext}{\textbf{Lesson Title}} &
\textcolor{headertext}{\textbf{Description}} &
\textcolor{headertext}{\textbf{Standard}} \\
\endhead

\multicolumn{5}{r}{\small\itshape Continued on next page\ldots} \\
\endfoot

\bottomrule
\endlastfoot

\toprule

% ── UNIT 1: Rigid Transformations and Congruence ─────────────────────────────
\multicolumn{5}{l}{\cellcolor{domainA}\textbf{Unit 1: Rigid Transformations and Congruence}} \\
\midrule

\rowcolor{row1}
1 & Unit 1: Rigid Transformations \& Congruence & Translations, Reflections \& Rotations &
Understand and apply translations, reflections, and rotations on the coordinate plane. Describe the effect of each rigid motion on coordinates and identify preserved properties (distance, angle measure). &
8.G.A.1 \\

\rowcolor{row2}
2 & Unit 1: Rigid Transformations \& Congruence & Sequences of Transformations &
Apply sequences of rigid motions to map one figure onto another. Understand that two figures are congruent if one can be obtained from the other by a sequence of rigid transformations; use this to identify congruent figures. &
8.G.A.2 \\

\rowcolor{row1}
3 & Unit 1: Rigid Transformations \& Congruence & Angle Relationships \& Parallel Lines &
Use informal arguments and transformation reasoning to establish facts about angles formed when parallel lines are cut by a transversal (alternate interior, corresponding angles) and the angle-sum of triangles. &
8.G.A.5 \\

\midrule
% ── UNIT 2: Dilations, Similarity, and Introducing Slope ─────────────────────
\multicolumn{5}{l}{\cellcolor{domainB}\textbf{Unit 2: Dilations, Similarity, and Introducing Slope}} \\
\midrule

\rowcolor{row1}
4 & Unit 2: Dilations, Similarity \& Introducing Slope & Dilations \& Scale Factor &
Understand dilations as transformations that scale figures from a center point by a scale factor. Identify the effect on coordinates and determine whether the image is larger, smaller, or congruent. &
8.G.A.3 / 8.G.A.4 \\

\rowcolor{row2}
5 & Unit 2: Dilations, Similarity \& Introducing Slope & Similar Figures &
Understand that two figures are similar if one can be obtained from the other by a sequence of rigid motions and a dilation. Use similarity to set up and solve proportions for unknown side lengths and angles. &
8.G.A.4 \\

\rowcolor{row1}
6 & Unit 2: Dilations, Similarity \& Introducing Slope & Introducing Slope via Similar Triangles &
Use similar triangles to explain why the slope $m$ is the same between any two distinct points on a non-vertical line in the coordinate plane. Derive the equation $y = mx$ for lines through the origin. &
8.EE.B.6 \\

\midrule
% ── UNIT 3: Linear Relationships ─────────────────────────────────────────────
\multicolumn{5}{l}{\cellcolor{domainC}\textbf{Unit 3: Linear Relationships}} \\
\midrule

\rowcolor{row1}
7 & Unit 3: Linear Relationships & Proportional Relationships \& Slope &
Graph proportional relationships; interpret the unit rate as the slope of the graph. Compare two proportional relationships represented in different ways (table, graph, equation, verbal). &
8.EE.B.5 \\

\rowcolor{row2}
8 & Unit 3: Linear Relationships & Slope-Intercept Form $y = mx + b$ &
Derive $y = mx + b$ as defining a linear function with slope $m$ and $y$-intercept $b$. Interpret slope and intercept in real-world contexts; graph and write equations given two points or a point and slope. &
8.EE.B.6 / 8.F.A.3 \\

\rowcolor{row1}
9 & Unit 3: Linear Relationships & Constructing \& Interpreting Linear Models &
Construct a function to model a linear relationship between two quantities. Determine rate of change and initial value from descriptions, tables, graphs, and equations; interpret in context. &
8.F.B.4 \\

\midrule
% ── UNIT 4: Linear Equations and Linear Systems ───────────────────────────────
\multicolumn{5}{l}{\cellcolor{domainD}\textbf{Unit 4: Linear Equations and Linear Systems}} \\
\midrule

\rowcolor{row1}
10 & Unit 4: Linear Equations \& Linear Systems & Solving Linear Equations in One Variable &
Solve linear equations in one variable including cases with one solution, infinitely many solutions, or no solution. Recognize when an equation has each type of solution set; justify reasoning. &
8.EE.C.7 \\

\rowcolor{row2}
11 & Unit 4: Linear Equations \& Linear Systems & Systems of Equations: Graphing &
Understand that a solution to a system of two linear equations is the point of intersection of their graphs. Solve systems graphically; interpret solutions in real-world contexts including no-solution and infinite-solution cases. &
8.EE.C.8a / 8.EE.C.8c \\

\rowcolor{row1}
12 & Unit 4: Linear Equations \& Linear Systems & Systems of Equations: Substitution \& Elimination &
Solve systems of two linear equations algebraically using substitution and elimination (linear combinations). Estimate solutions by graphing and verify algebraically. &
8.EE.C.8b / 8.EE.C.8c \\

\midrule
% ── UNIT 5: Functions and Volume ─────────────────────────────────────────────
\multicolumn{5}{l}{\cellcolor{domainE}\textbf{Unit 5: Functions and Volume}} \\
\midrule

\rowcolor{row1}
13 & Unit 5: Functions \& Volume & Understanding Functions &
Understand that a function assigns exactly one output to each input. Represent functions using tables, graphs, and equations; determine whether a relation is a function. Distinguish linear from nonlinear functions. &
8.F.A.1 / 8.F.A.2 / 8.F.A.3 \\

\rowcolor{row2}
14 & Unit 5: Functions \& Volume & Qualitative Graphs \& Functional Relationships &
Describe qualitatively the functional relationship between two quantities by analyzing a graph (e.g., where it is increasing, decreasing, linear, or nonlinear). Sketch a graph from a verbal description. &
8.F.B.5 \\

\rowcolor{row1}
15 & Unit 5: Functions \& Volume & Volume of Cylinders, Cones \& Spheres &
Know and apply the formulas for volumes of cylinders ($V{=}\pi r^2 h$), cones ($V{=}\tfrac{1}{3}\pi r^2 h$), and spheres ($V{=}\tfrac{4}{3}\pi r^3$). Solve real-world and mathematical problems. &
8.G.C.9 \\

\midrule
% ── UNIT 6: Associations in Data ─────────────────────────────────────────────
\multicolumn{5}{l}{\cellcolor{domainF}\textbf{Unit 6: Associations in Data}} \\
\midrule

\rowcolor{row1}
16 & Unit 6: Associations in Data & Scatter Plots \& Patterns of Association &
Construct and interpret scatter plots for bivariate measurement data. Describe patterns of association including clustering, outliers, positive/negative correlation, and linear vs.\ nonlinear trends. &
8.SP.A.1 \\

\rowcolor{row2}
17 & Unit 6: Associations in Data & Lines of Best Fit \& Predictions &
Informally fit a straight line to data suggesting a linear association. Interpret slope and intercept in context; use the line of best fit to make predictions and assess model fit. &
8.SP.A.2 / 8.SP.A.3 \\

\rowcolor{row1}
18 & Unit 6: Associations in Data & Two-Way Frequency Tables &
Construct and interpret two-way tables summarizing bivariate categorical data from a sample. Use relative frequencies to investigate possible associations between two categorical variables. &
8.SP.A.4 \\

\midrule
% ── UNIT 7: Exponents and Scientific Notation ────────────────────────────────
\multicolumn{5}{l}{\cellcolor{domainG}\textbf{Unit 7: Exponents and Scientific Notation}} \\
\midrule

\rowcolor{row1}
19 & Unit 7: Exponents \& Scientific Notation & Properties of Integer Exponents &
Know and apply properties of integer exponents (product, quotient, power, zero, and negative exponent rules) to generate equivalent numerical expressions (e.g., $3^2 \cdot 3^{-5} = 3^{-3} = \tfrac{1}{27}$). &
8.EE.A.1 \\

\rowcolor{row2}
20 & Unit 7: Exponents \& Scientific Notation & Scientific Notation \& Orders of Magnitude &
Express very large and very small quantities in scientific notation. Choose units of appropriate size; compare magnitudes and perform operations (add, subtract, multiply, divide) with numbers in scientific notation. &
8.EE.A.3 / 8.EE.A.4 \\

\midrule
% ── UNIT 8: Pythagorean Theorem and Irrational Numbers ───────────────────────
\multicolumn{5}{l}{\cellcolor{domainH}\textbf{Unit 8: Pythagorean Theorem and Irrational Numbers}} \\
\midrule

\rowcolor{row1}
21 & Unit 8: Pythagorean Theorem \& Irrational Numbers & Square \& Cube Roots &
Use square root and cube root symbols to represent solutions to $x^2 = p$ and $x^3 = p$. Evaluate roots of perfect squares and cubes; recognize $\sqrt{2}$ as irrational and connect to the Pythagorean Theorem. &
8.EE.A.2 \\

\rowcolor{row2}
22 & Unit 8: Pythagorean Theorem \& Irrational Numbers & The Pythagorean Theorem &
Explain a proof of the Pythagorean Theorem and its converse. Apply to determine unknown side lengths and classify triangles (right, acute, obtuse) in real-world and mathematical problems. &
8.G.B.7 \\

\rowcolor{row1}
23 & Unit 8: Pythagorean Theorem \& Irrational Numbers & Distance in the Coordinate Plane &
Apply the Pythagorean Theorem to find the distance between two points in the coordinate plane. Derive and use the distance formula in real-world contexts. &
8.G.B.8 \\

\rowcolor{row2}
24 & Unit 8: Pythagorean Theorem \& Irrational Numbers & Rational \& Irrational Numbers &
Distinguish rational from irrational numbers. Understand that irrational numbers have non-repeating, non-terminating decimal expansions. Use rational approximations of irrationals to compare size and locate on a number line. &
8.NS.A.1 / 8.NS.A.2 \\

\midrule
% ── UNIT 9: Capstone ─────────────────────────────────────────────────────────
\multicolumn{5}{l}{\cellcolor{domainI}\textbf{Unit 9: Capstone --- When Am I Ever Going to Use This?}} \\
\midrule

\rowcolor{row1}
25 & Unit 9: Capstone & Cumulative Real-World Applications &
Solve complex, multi-step real-world problems integrating all 8th grade domains: transformations, linear equations, systems, functions, volume, data analysis, exponents, and the Pythagorean Theorem. Communicate mathematical reasoning clearly and justify solutions. &
8.MP.1--8.MP.8 \\

\end{longtable}

\end{document}